\subsection{Normal Operator Consequences}

All inner-product spaces in this subsection are finite-dimensional over the stated field.

\begin{perspective} \label{per:normal-operators-coordinate-free-diagonal} \hlt{Normal Means Diagonal in Orthogonal Coordinates}\\
Theorem \ref{thm:complex-spectral-theorem} identifies complex normality with orthonormal diagonal form.\\
Thus normality is the coordinate-free form of orthogonal diagonalizability.\\
Over $\mathbb{R}$, invariant planes replace the missing nonreal eigenvectors.
\end{perspective}

\begin{proposition} \label{prop:normal-polynomial-kernel-range-properties} \hlt{Basic Spectral Properties of Normal Operators}\\
Let $\mathcal{V}$ be a complex inner-product space, $T\in\LL(\mathcal{V})$ be normal, and $p\in\mathbb{C}[z]$.
\begin{enumerate}[label=(\roman*)]
\setlength{\itemsep}{0pt}
\item $p(T)$ is normal and commutes with $T$.
\item $T^*$ is normal and $(\forall\lambda\in\mathbb{C})\ E_{\lambda}(T)=E_{\overline{\lambda}}(T^*)$.
\item $\ker T=\ker T^*$ and $\operatorname{range}(T)=\operatorname{range}(T^*)$.
\end{enumerate}
\begin{proof}
Set $n\defeq\dim(\mathcal{V})$. By Theorem \ref{thm:complex-spectral-theorem}(i),(ii), fix an orthonormal basis $(\mathbf{e}_j)_{j=1}^{n}$ satisfying $(\forall j\in J_n)\,T\mathbf{e}_j=\mu_j\mathbf{e}_j$.
\begin{enumerate}[label=(\roman*)]
\setlength{\itemsep}{0pt}
\item $(\forall j\in J_n)\,p(T)\mathbf{e}_j=p(\mu_j)\mathbf{e}_j$ by polynomial evaluation. The operator $p(T)$ is normal by Theorem \ref{thm:complex-spectral-theorem}(iii),(i).\\
The identity $Tp(T)=p(T)T$ holds by Proposition \ref{prop:operator-polynomial-products}.
\item $(\forall j\in J_n)\,T^*\mathbf{e}_j=\overline{\mu_j}\mathbf{e}_j$ by Theorem \ref{thm:normal-operator-adjoint-eigenvectors}. Thus $T^*$ is normal by Theorem \ref{thm:complex-spectral-theorem}(iii),(i).\\
For $\lambda\in\mathbb{C}$, $E_{\lambda}(T)\subseteq E_{\overline{\lambda}}(T^*)$ by Theorem \ref{thm:normal-operator-adjoint-eigenvectors}.\\
The reverse inclusion holds by the same theorem for $T^*$ and $(T^*)^*=T$ by Proposition \ref{prop:basic-adjoint-identities}(iii); hence equality.
\item $\ker T=E_0(T)=E_0(T^*)=\ker T^*$ by (ii).\\
$\operatorname{range}(T)=\operatorname{span}\{\mathbf{e}_j:\mu_j\neq0\}=\operatorname{range}(T^*)$ by the common eigenbasis and $\mu_j\neq0\Longleftrightarrow\overline{\mu_j}\neq0$.
\end{enumerate}
\hlt{Structural Commentary}: The spectral theorem transfers polynomial, adjoint, kernel, and range claims to one diagonal basis.\\
\hlt{Key Insight}: Adjunction conjugates eigenvalues while preserving the zero eigenspace.
\end{proof}
\end{proposition}

\begin{lemma} \label{lem:adjoints-of-finitely-many-normal-matrices-polynomial} \hlt{Adjoints of Normal Matrices as One Polynomial}\\
Let $k\in\mathbb{Z}_{\geq1}$ and $(n_j)_{j=1}^{k}\in\mathbb{Z}_{\geq1}^{k}$. Let $(A_j)_{j=1}^{k}$ satisfy $(\forall j\in J_k)\,[A_j\in\operatorname{Mat}(n_j,n_j,\mathbb{C})\land A_j^*A_j=A_jA_j^*]$.\\
Then $(\exists p\in\mathbb{C}[z])(\forall j\in J_k)\,p(A_j)=A_j^*$.
\begin{proof}
\begin{enumerate}[label=(\roman*)]
\setlength{\itemsep}{0pt}
\item \hlt{Spectral Set}: Set $\Omega\defeq\bigcup\limits_{j=1}^{k}\operatorname{spec}(A_j)$. By Theorem \ref{thm:complex-operator-has-eigenvalue} and Corollary \ref{cor:number-of-distinct-eigenvalues}, $(\forall j\in J_k)\,1\leq\#\operatorname{spec}(A_j)\leq n_j$; hence $0<\#\Omega<\infty$.
\item \hlt{Interpolation}: Define $p(z)\defeq\sum\limits_{\lambda\in\Omega}\overline{\lambda}\prod\limits_{\nu\in\Omega\setminus\{\lambda\}}\frac{z-\nu}{\lambda-\nu}$. Then $(\forall\lambda\in\Omega)\,p(\lambda)=\overline{\lambda}$ by direct evaluation of the Lagrange factors.
\item \hlt{Spectral Evaluation}: Fix $j\in J_k$. The matrix $A_j$ induces a normal operator on $\mathbb{C}^{n_j}$ by Theorem \ref{thm:matrix-of-adjoint-is-conjugate-transpose} and Proposition \ref{prop:composition-and-matrix-multiplication}.\\
Fix $A_j=\sum\limits_{\mu\in\operatorname{spec}(A_j)}\mu P_{j,\mu}$ by Theorem \ref{thm:spectral-resolution-normal-operator}(i). By Corollary \ref{cor:normal-spectral-resolution-functional-calculus} and the interpolation identity, $p(A_j)=\sum\limits_{\mu\in\operatorname{spec}(A_j)}\overline{\mu}P_{j,\mu}$.\\
By Proposition \ref{prop:basic-adjoint-identities}(i) and Theorem \ref{thm:orthogonal-projection-self-adjoint-idempotent}, $A_j^*=\sum\limits_{\mu\in\operatorname{spec}(A_j)}\overline{\mu}P_{j,\mu}$. Thus $(\forall j\in J_k)\,p(A_j)=A_j^*$.
\end{enumerate}
\hlt{Structural Commentary}: Spectral union and Lagrange interpolation precede evaluation on each normal matrix.\\
\hlt{Key Insight}: A finite spectral union admits one conjugating polynomial.
\end{proof}
\end{lemma}

\begin{theorem} \label{thm:fuglede-putnam-normal-intertwining} \hlt{Fuglede--Putnam Theorem}\\
Let $m,n\geq1$, and let $A\in\operatorname{Mat}(m,m,\mathbb{C})$ and $B\in\operatorname{Mat}(n,n,\mathbb{C})$ be normal. Let $X\in\operatorname{Mat}(m,n,\mathbb{C})$. Then $AX=XB\Longleftrightarrow A^*X=XB^*$.
\begin{proof}
\begin{enumerate}[label=\roman*.]
\setlength{\itemsep}{0pt}
\item[$(\Rightarrow)$] Assume $AX=XB$. Then $(\forall r\in\mathbb{Z}_{\geq0})\,A^rX=XB^r$ by induction, so $(\forall p\in\mathbb{C}[z])\,p(A)X=Xp(B)$ by linearity.\\
Fix $p$ with $p(A)=A^*$ and $p(B)=B^*$ by Lemma \ref{lem:adjoints-of-finitely-many-normal-matrices-polynomial}. Then $A^*X=p(A)X=Xp(B)=XB^*$.
\item[$(\Leftarrow)$] Assume $A^*X=XB^*$. The matrices $A^*$ and $B^*$ are normal by $(A^*)^*A^*=AA^*=A^*A=A^*(A^*)^*$ and its analogue for $B^*$.\\
Thus $(\forall q\in\mathbb{C}[z])\,q(A^*)X=Xq(B^*)$ by induction on powers and linearity.\\
Fix $q$ with $q(A^*)=A$ and $q(B^*)=B$ by Lemma \ref{lem:adjoints-of-finitely-many-normal-matrices-polynomial} and Proposition \ref{prop:basic-adjoint-identities}(iii). Then $AX=q(A^*)X=Xq(B^*)=XB$.
\end{enumerate}
\hlt{Structural Commentary}: Polynomial propagation and common interpolation form symmetric directional arguments.\\
\hlt{Key Insight}: Normal adjoints are polynomial functions of their matrices.
\end{proof}
\end{theorem}

\begin{computationmethod} \label{meth:finding-spectral-decomposition-normal-matrix} \hlt{Finding a Spectral Decomposition of a Normal Matrix}\\
Let $n\geq1$ and let $A\in\operatorname{Mat}(n,n,\mathbb{C})$ be normal. Its eigenspaces span $\mathbb{C}^n$ by Theorem \ref{thm:complex-spectral-theorem}(i),(ii).\\
Distinct eigenspaces are orthogonal by Theorem \ref{thm:normal-eigenvectors-distinct-eigenvalues-orthogonal}. The output is a unitary diagonalization $A=U\Lambda U^*$.
\begin{enumerate}[label=\roman*.]
\setlength{\itemsep}{0pt}
\item Find the distinct eigenvalues $\mu_1,\ldots,\mu_d$.
\item Compute $(E_{\mu_j}(A))_{j=1}^{d}$ by solving $(A-\mu_jI_n)\mathbf{x}=0$ for $j\in J_d$.
\item Orthonormalize bases of $(E_{\mu_j}(A))_{j=1}^{d}$.
\item Concatenate these bases as the columns of $U$, and set $\Lambda\defeq U^*AU$.
\item Verify $U^*U=UU^*=I_n$ and $AU=U\Lambda$. Check that the diagonal of $\Lambda$ repeats $\mu_j$ exactly $\dim E_{\mu_j}(A)$ times for $j\in J_d$.
\end{enumerate}
\end{computationmethod}

\begin{theorem} \label{thm:normal-spectral-classification-by-eigenvalues} \hlt{Spectral Classification of Normal Operators}\\
Let $\mathcal{V}$ be a finite-dimensional complex inner-product space and let $T\in\LL(\mathcal{V})$ be normal.
\begin{enumerate}[label=(\roman*)]
\setlength{\itemsep}{0pt}
\item $T$ is unitary if and only if $(\forall\lambda\in\operatorname{spec}(T))\,\abs{\lambda}=1$.
\item $T$ is self-adjoint if and only if $\operatorname{spec}(T)\subseteq\mathbb{R}$.
\item $T$ is an orthogonal projection if and only if $\operatorname{spec}(T)\subseteq\{0,1\}$.
\item $T^*=-T$ if and only if $\operatorname{spec}(T)\subseteq i\mathbb{R}$.
\end{enumerate}
\begin{proof}
If $\mathcal{V}=\{\mathbf{0}\}$, its unique operator satisfies all four operator conditions.\\
Its empty spectrum satisfies all four spectral conditions. Assume $\mathcal{V}\neq\{\mathbf{0}\}$.\\
List the distinct eigenvalues as $\mu_1,\ldots,\mu_d$, and fix $T=\sum\limits_{j=1}^{d}\mu_jP_j$ by Theorem \ref{thm:spectral-resolution-normal-operator}(i).\\
The identity $I_{\mathcal{V}}=\sum\limits_{j=1}^{d}P_j$ holds by Theorem \ref{thm:spectral-resolution-normal-operator}(i). The formula $T^*=\sum\limits_{j=1}^{d}\overline{\mu_j}P_j$ holds by Proposition \ref{prop:basic-adjoint-identities}(i) and Theorem \ref{thm:orthogonal-projection-self-adjoint-idempotent}.\\
The formulas $T^*T=\sum\limits_{j=1}^{d}\abs{\mu_j}^2P_j$ and $T^2=\sum\limits_{j=1}^{d}\mu_j^2P_j$ hold by direct multiplication.\\
To verify linear independence, suppose $\sum\limits_{j=1}^{d}\alpha_jP_j=0$ and fix $k\in J_d$.\\
Fix $\mathbf{0}\neq\mathbf{x}\in E_{\mu_k}(T)=\operatorname{range}(P_k)$ by $\mu_k\in\operatorname{spec}(T)$ and Theorem \ref{thm:spectral-resolution-normal-operator}(ii).\\
Then $0=\sum\limits_{j=1}^{d}\alpha_jP_j\mathbf{x}=\alpha_k\mathbf{x}$, so $\alpha_k=0$. Thus $(P_j)_{j=1}^{d}$ is linearly independent.
\begin{enumerate}[label=(\roman*)]
\setlength{\itemsep}{0pt}
\item[$(i)$] ($\Rightarrow$) If $T$ is unitary, then $T^*T=I_{\mathcal{V}}$ by unitarity.\\
Hence $(\forall j\in J_d)\,\abs{\mu_j}^2=1$ by linear independence of the $P_j$, and therefore $(\forall j\in J_d)\,\abs{\mu_j}=1$.
\item[$(i)$] ($\Leftarrow$) If $(\forall j\in J_d)\,\abs{\mu_j}=1$, then $T^*T=TT^*=\sum\limits_{j=1}^{d}P_j=I_{\mathcal{V}}$. Thus $T$ is unitary by the definition of unitarity.
\item[$(ii)$] ($\Rightarrow$) If $T=T^*$, then $\sum\limits_{j=1}^{d}(\mu_j-\overline{\mu_j})P_j=0$. Hence $(\forall j\in J_d)\,\mu_j=\overline{\mu_j}$ by linear independence of the $P_j$, so $\operatorname{spec}(T)\subseteq\mathbb{R}$.
\item[$(ii)$] ($\Leftarrow$) If $\operatorname{spec}(T)\subseteq\mathbb{R}$, then $T^*=\sum\limits_{j=1}^{d}\overline{\mu_j}P_j=\sum\limits_{j=1}^{d}\mu_jP_j=T$.
\item[$(iii)$] ($\Rightarrow$) If $T$ is an orthogonal projection, then $T^*=T$ and $T^2=T$ by Theorem \ref{thm:orthogonal-projection-self-adjoint-idempotent}.\\
Hence $(\forall j\in J_d)\,\mu_j=\overline{\mu_j}$ and $\mu_j^2=\mu_j$ by linear independence of the $P_j$. Therefore $\operatorname{spec}(T)\subseteq\{0,1\}$.
\item[$(iii)$] ($\Leftarrow$) If $\operatorname{spec}(T)\subseteq\{0,1\}$, then $(\forall j\in J_d)\,\overline{\mu_j}=\mu_j^2=\mu_j$. Hence $T^*=T^2=T$ by the spectral formulas.\\
Thus $T$ is an orthogonal projection by Theorem \ref{thm:orthogonal-projection-self-adjoint-idempotent}.
\item[$(iv)$] ($\Rightarrow$) If $T^*=-T$, then $\sum\limits_{j=1}^{d}(\overline{\mu_j}+\mu_j)P_j=0$. Hence $(\forall j\in J_d)\,\overline{\mu_j}=-\mu_j$ by linear independence of the $P_j$, so $\operatorname{spec}(T)\subseteq i\mathbb{R}$.
\item[$(iv)$] ($\Leftarrow$) If $\operatorname{spec}(T)\subseteq i\mathbb{R}$, then $(\forall j\in J_d)\,\overline{\mu_j}=-\mu_j$. Consequently, $T^*=\sum\limits_{j=1}^{d}\overline{\mu_j}P_j=-\sum\limits_{j=1}^{d}\mu_jP_j=-T$.
\end{enumerate}
\hlt{Structural Commentary}: One spectral resolution converts four operator identities into coefficient comparisons.\\
\hlt{Key Insight}: Orthogonal spectral projections are linearly independent.
\end{proof}
\end{theorem}

\begin{proposition} \label{prop:cartesian-decomposition-normal-commuting-parts} \hlt{Cartesian Decomposition and Normality}\\
Let $\mathcal{V}$ be a finite-dimensional complex inner-product space and let $T\in\LL(\mathcal{V})$.
\begin{enumerate}[label=(\roman*)]
\setlength{\itemsep}{0pt}
\item The operator $T$ has a unique decomposition $T=H+iK$ with self-adjoint $H,K\in\LL(\mathcal{V})$.\\
Its terms are $H=\frac{1}{2}(T+T^*)$ and $K=\frac{1}{2i}(T-T^*)$.
\item For these operators, $T$ is normal if and only if $HK=KH$.
\end{enumerate}
\begin{proof}
\begin{enumerate}[label=(\roman*)]
\setlength{\itemsep}{0pt}
\item[$(i)$] Set $H\defeq\frac{1}{2}(T+T^*)$ and $K\defeq\frac{1}{2i}(T-T^*)$. Then $H^*=H$, $K^*=K$, and $T=H+iK$ by Proposition \ref{prop:basic-adjoint-identities}(i),(iii).\\
If $T=X+iY$ for self-adjoint $X,Y$, then $T^*=X-iY$ by Proposition \ref{prop:basic-adjoint-identities}(i).\\
Consequently, $X=\frac{1}{2}(T+T^*)=H$ and $Y=\frac{1}{2i}(T-T^*)=K$.
\item[$(ii)$] ($\Rightarrow$) Suppose $T$ is normal.\\
Then $HK=KH$ by the identities $T^*T=H^2+K^2+i(HK-KH)$, $TT^*=H^2+K^2-i(HK-KH)$, and $T^*T=TT^*$.
\item[$(ii)$] ($\Leftarrow$) If $HK=KH$, the same product identities satisfy $T^*T=TT^*$, so $T$ is normal.
\end{enumerate}
\hlt{Structural Commentary}: Adjoint symmetrization precedes a direct product comparison.\\
\hlt{Key Insight}: The normality defect is $2i(HK-KH)$.
\end{proof}
\end{proposition}

Figure \ref{fig:cartesian-normal-commutation} summarizes Proposition \ref{prop:cartesian-decomposition-normal-commuting-parts}.

\begin{illustrationfigure}[H]
  \centering
  \begin{tikzpicture}[
  >=Stealth,
  every node/.style={font=\small},
  box/.style={draw=gray!55, rounded corners=3pt, fill=gray!6, align=center, inner xsep=7pt, inner ysep=4.5pt, minimum height=0.9cm},
  result/.style={box, fill=green!7, minimum width=4.05cm},
  arrow/.style={->, line width=0.85pt, blue!70!black}
]
  \node[box, minimum width=1.4cm] (T) at (0,0) {$T$};
  \node[box, minimum width=3.75cm] (parts) at (2.95,0)
    {$H=\frac12(T+T^*)$\\[-1pt]
     $K=\frac{1}{2i}(T-T^*)$\\[-1pt]
     $T=H+iK$};
  \node[result] (normal) at (7.35,0)
    {$T\ \text{normal}\quad\Longleftrightarrow\quad HK=KH$};

  \draw[arrow] (T) -- (parts);
  \draw[arrow] (parts) -- (normal);
\end{tikzpicture}

  \caption{Cartesian decomposition: $T$ is normal exactly when $HK=KH$.}
  \label{fig:cartesian-normal-commutation}
\end{illustrationfigure}

\begin{proposition} \label{prop:block-upper-triangular-normal-collapse} \hlt{Normal Block Upper-Triangular Matrices Collapse}\\
Let $\mathbb{F}\in\{\mathbb{R},\mathbb{C}\}$ and $r,s\geq1$. Let $(B,C,X)\in\operatorname{Mat}(r,r,\mathbb{F})\times\operatorname{Mat}(s,s,\mathbb{F})\times\operatorname{Mat}(r,s,\mathbb{F})$.\\
Set $A\defeq\left(\begin{smallmatrix}B&X\\0&C\end{smallmatrix}\right)\in\operatorname{Mat}(r+s,r+s,\mathbb{F})$. Then $A$ is normal if and only if $X=0$ and both $B$ and $C$ are normal.
\begin{proof}
\begin{enumerate}[label=\roman*.]
\setlength{\itemsep}{0pt}
\item[$(\Rightarrow)$] Suppose $A$ is normal; regard all matrices as complex when $\mathbb{F}=\mathbb{R}$.\\
By Theorem \ref{thm:schur-triangularization}, fix unitary $U_B,U_C$ such that $U_B^*BU_B$ and $U_C^*CU_C$ are upper triangular.\\
Set $Q\defeq U_B\oplus U_C$, which is unitary by Proposition \ref{prop:algebra-of-unitary-matrices}(ii). Then $Q^*Q=QQ^*=I_{r+s}$ by Theorem \ref{thm:unitary-matrix-equivalent-conditions}(i),(iii).\\
The matrix $Q^*AQ=\left(\begin{smallmatrix}U_B^*BU_B&U_B^*XU_C\\0&U_C^*CU_C\end{smallmatrix}\right)$ is upper triangular.\\
The matrix $Q^*AQ$ is normal by $(Q^*AQ)^*(Q^*AQ)=Q^*A^*AQ=Q^*AA^*Q=(Q^*AQ)(Q^*AQ)^*$.\\
Thus $Q^*AQ$ is diagonal by Lemma \ref{lem:upper-triangular-normal-matrix-is-diagonal}. Hence $U_B^*XU_C=0$ and $X=U_B(U_B^*XU_C)U_C^*=0$.\\
Now $A=B\oplus C$, so $B^*B=BB^*$ and $C^*C=CC^*$ by blockwise normality.
\item[$(\Leftarrow)$] If $X=0$ and $B,C$ are normal, then $A=B\oplus C$ and $A^*A=AA^*$ by blockwise multiplication.
\end{enumerate}
\hlt{Structural Commentary}: Blockwise Schur reduction converts normality into the triangular-collapse lemma.\\
\hlt{Key Insight}: A normal triangular coupling cannot survive diagonalization.
\end{proof}
\end{proposition}

Figure \ref{fig:block-upper-triangular-normal-collapse} summarizes Proposition \ref{prop:block-upper-triangular-normal-collapse}. Block-diagonal Schur coordinates reduce the claim to the earlier triangular normal-matrix lemma.

\begin{illustrationfigure}[H]
  \centering
  \begin{tikzpicture}[
  >=Stealth,
  every node/.style={font=\small},
  block/.style={draw=gray!55, rounded corners=3pt, inner xsep=6pt, inner ysep=5pt, align=center},
  arrow/.style={->, line width=0.85pt, blue!70!black}
]
  \node[block] (tri) at (0,0)
    {$A=\begin{bmatrix}B&X\\0&C\end{bmatrix}$\\normal};
  \node[block] (schur) at (3.55,0)
    {$T_B=U_B^*BU_B,\quad T_C=U_C^*CU_C$\\
     $Q^*AQ=\begin{bmatrix}T_B&U_B^*XU_C\\0&T_C\end{bmatrix}$\\
     upper triangular and normal};
  \node[block] (diag) at (7.62,0)
    {$Q^*AQ$ diagonal\\$X=0$, $B,C$ normal};

  \draw[arrow] (tri) -- (schur);
  \draw[arrow] (schur) -- (diag);
\end{tikzpicture}

  \caption{Normality removes the upper-right block in a block upper-triangular matrix.}
  \label{fig:block-upper-triangular-normal-collapse}
\end{illustrationfigure}

\begin{theorem} \label{thm:normal-invariant-subspaces-reduce} \hlt{Invariant Subspaces of Normal Matrices Reduce}\\
Let $n\geq1$, $\mathbb{F}\in\{\mathbb{R},\mathbb{C}\}$, and $A\in\operatorname{Mat}(n,n,\mathbb{F})$ be normal. Let $\mathcal{U}\leq\mathbb{F}^n$ be invariant under $A$, and let $P$ be its orthogonal projection.
\begin{enumerate}[label=(\roman*)]
\setlength{\itemsep}{0pt}
\item The subspace $\mathcal{U}$ is invariant under $A^*$.
\item The identity $AP=PA$ holds.
\end{enumerate}
\begin{proof}
If $\mathcal{U}=\{\mathbf{0}\}$, then $A^*\mathcal{U}\subseteq\mathcal{U}$ and $P=0$, so both claims hold. If $\mathcal{U}=\mathbb{F}^n$, then $A^*\mathcal{U}\subseteq\mathcal{U}$ and $P=I_n$, so both claims hold.\\
Assume $r\defeq\dim(\mathcal{U})\in\{1,\ldots,n-1\}$. Choose orthonormal bases of $\mathcal{U}$ and $\mathcal{U}^{\perp}$ by Corollary \ref{cor:finite-dimensional-inner-product-space-has-orthonormal-basis}.\\
Let $Q=[Q_1\ Q_2]$ be the unitary matrix obtained by concatenating these bases.\\
The identity $Q_2^*AQ_1=0$ holds by invariance of $\mathcal{U}$, so $Q^*AQ=\left(\begin{smallmatrix}B&X\\0&C\end{smallmatrix}\right)$.\\
The matrix $Q^*AQ$ is normal by $(Q^*AQ)^*(Q^*AQ)=Q^*A^*AQ$ and $(Q^*AQ)(Q^*AQ)^*=Q^*AA^*Q$. Hence $X=0$ by Proposition \ref{prop:block-upper-triangular-normal-collapse}.
\begin{enumerate}[label=(\roman*)]
\setlength{\itemsep}{0pt}
\item The inclusion $A^*\mathcal{U}\subseteq\mathcal{U}$ holds by $Q^*A^*Q=(Q^*AQ)^*=B^*\oplus C^*$.
\item The identity $Q^*PQ=I_r\oplus0$ holds by Proposition \ref{prop:orthogonal-projection-properties}(ii)--(iii).\\
The matrices $Q^*AQ$ and $Q^*PQ$ commute by blockwise multiplication; hence $AP=PA$ by conjugation with $Q$.
\end{enumerate}
\hlt{Structural Commentary}: An adapted unitary basis reduces both conclusions to block multiplication.\\
\hlt{Key Insight}: Normality eliminates the off-diagonal coupling block.
\end{proof}
\end{theorem}

\begin{corollary} \label{cor:normal-operator-invariant-subspaces-reduce} \hlt{Invariant Subspaces of Normal Operators Reduce}\\
Let $\mathbb{F}\in\{\mathbb{R},\mathbb{C}\}$, let $\mathcal{V}$ be an inner-product space over $\mathbb{F}$, and let $T\in\LL(\mathcal{V})$ be normal. Let $\mathcal{U}\leq\mathcal{V}$ be invariant under $T$.
\begin{enumerate}[label=(\roman*)]
\setlength{\itemsep}{0pt}
\item $\mathcal{U}^{\perp}$ is invariant under $T$;
\item $\mathcal{U}$ is invariant under $T^*$;
\item $(T|_{\mathcal{U}})^*=(T^*)|_{\mathcal{U}}$;
\item $T|_{\mathcal{U}}$ and $T|_{\mathcal{U}^{\perp}}$ are normal.
\end{enumerate}
\begin{proof}
If $\mathcal{U}=\{\mathbf{0}\}$ or $\mathcal{U}=\mathcal{V}$, all four claims hold by $T\mathbf{0}=T^*\mathbf{0}=\mathbf{0}$ and normality of the restrictions.\\
Assume $0<r\defeq\dim(\mathcal{U})<n\defeq\dim(\mathcal{V})$. Choose an orthonormal basis $\beta$ adapted to $\mathcal{U}\oplus\mathcal{U}^{\perp}$ by Corollary \ref{cor:finite-dimensional-inner-product-space-has-orthonormal-basis} and Theorem \ref{thm:orthogonal-direct-sum-decomposition}.\\
Set $A\defeq[T]_{\beta\leftarrow\beta}$ and $R\defeq I_r\oplus0_{n-r}$. The matrix $A$ is normal by Theorem \ref{thm:matrix-of-adjoint-is-conjugate-transpose} and Proposition \ref{prop:composition-and-matrix-multiplication}.\\
The coordinate subspace $\mathbb{F}^r\oplus\{\mathbf{0}\}$ is invariant under $A$, and $R$ is its orthogonal projection.\\
Hence $AR=RA$ by Theorem \ref{thm:normal-invariant-subspaces-reduce}(ii). The identity $A=B\oplus C$ holds by direct block multiplication with $R$.
\begin{enumerate}[label=(\roman*)]
\setlength{\itemsep}{0pt}
\item The inclusion $T(\mathcal{U}^{\perp})\subseteq\mathcal{U}^{\perp}$ holds by $A=B\oplus C$.
\item The identity $[T^*]_{\beta\leftarrow\beta}=A^*=B^*\oplus C^*$ holds by Theorem \ref{thm:matrix-of-adjoint-is-conjugate-transpose}. Hence $\mathcal{U}$ is invariant under $T^*$.
\item The subspaces $\mathcal{U}$ and $\mathcal{U}^{\perp}$ are invariant under $T$ by the hypothesis and part (i).\\
The identity $\mathcal{V}=\mathcal{U}\oplus\mathcal{U}^{\perp}$ holds by Theorem \ref{thm:orthogonal-direct-sum-decomposition}. Thus $(\mathcal{U},\mathcal{U}^{\perp})$ reduces $T$, and $(T|_{\mathcal{U}})^*=(T^*)|_{\mathcal{U}}$ by Theorem \ref{thm:adjoint-invariant-subspace-orthogonal-complement}(iii).
\item The identity $A^*A=AA^*$ is equivalent blockwise to $B^*B=BB^*$ and $C^*C=CC^*$.\\
By Theorem \ref{thm:matrix-of-adjoint-is-conjugate-transpose} and Proposition \ref{prop:composition-and-matrix-multiplication}, these are the adjoint-product identities for both restrictions.\\
Hence both restrictions are normal by Theorem \ref{thm:linear-maps-and-matrices-same-data}(i).
\end{enumerate}
\hlt{Structural Commentary}: Matrix reduction transfers four operator claims to two diagonal blocks.\\
\hlt{Key Insight}: Commutation with the projection is the reducing criterion.
\end{proof}
\end{corollary}

\begin{theorem} \label{thm:simultaneous-unitary-diagonalization-normal-family} \hlt{Simultaneous Unitary Diagonalization of Commuting Normal Operators}\\
Let $\mathcal{V}$ be a complex inner-product space, and let $\emptyset\neq\mathcal{F}\subseteq\LL(\mathcal{V})$ consist of normal operators.\\
Then $\mathcal{F}$ is pairwise commuting if and only if $\mathcal{V}$ has an orthonormal basis of common eigenvectors for $\mathcal{F}$.
\begin{proof}
\begin{enumerate}[label=\roman*.]
\setlength{\itemsep}{0pt}
\item[$(\Rightarrow)$] Assume $\mathcal{F}$ is pairwise commuting and induct on $\dim(\mathcal{V})$.\\
If $(\forall A\in\mathcal{F})(\exists c_A\in\mathbb{C})\,A=c_AI_{\mathcal{V}}$, fix any orthonormal basis by Corollary \ref{cor:finite-dimensional-inner-product-space-has-orthonormal-basis}.\\
It is common by the scalar condition, including the empty basis when $\mathcal{V}=\{\mathbf{0}\}$.\\
Otherwise, fix a nonscalar $A_0\in\mathcal{F}$ and its spectral resolution $A_0=\sum\limits_{j=1}^{d}\mu_jP_j$ by Theorem \ref{thm:spectral-resolution-normal-operator}(i).\\
The inequality $d\geq2$ holds by nonscalarity. $(\forall j\in J_d)\,0<\dim(\operatorname{range}(P_j))<\dim(\mathcal{V})$ by Theorem \ref{thm:spectral-resolution-normal-operator}(i)--(ii) and $d\geq2$.\\
Each $P_j$ is a polynomial in $A_0$ by Corollary \ref{cor:normal-spectral-resolution-functional-calculus}(ii).\\
Fix $B\in\mathcal{F}$. Then $(\forall p\in\mathbb{C}[z])\,Bp(A_0)=p(A_0)B$ by $BA_0=A_0B$, induction on powers, and linearity.\\
Hence $(\forall j\in J_d)\,BP_j=P_jB$. Fix $j\in J_d$ and $\mathbf{x}\in\operatorname{range}(P_j)$. Then $P_jB\mathbf{x}=BP_j\mathbf{x}=B\mathbf{x}$, so $B\mathbf{x}\in\operatorname{range}(P_j)$.\\
Thus $B|_{\operatorname{range}(P_j)}$ is normal by Corollary \ref{cor:normal-operator-invariant-subspaces-reduce}(iv). The restricted family on each range is pairwise commuting and normal.\\
Fix a common orthonormal eigenbasis on each range by the induction hypothesis.\\
The identity $\mathcal{V}=\bigoplus\limits_{j=1}^{d}{}^{\perp}\operatorname{range}(P_j)$ holds by Theorem \ref{thm:spectral-resolution-normal-operator}(i)--(ii).\\
Concatenating the restricted bases produces a common orthonormal eigenbasis of $\mathcal{V}$.
\item[$(\Leftarrow)$] If $\mathcal{V}=\{\mathbf{0}\}$, its unique operator commutes with itself.\\
Otherwise, every $A\in\mathcal{F}$ has a diagonal matrix in the common orthonormal eigenbasis.\\
Hence the operators commute by Proposition \ref{prop:composition-and-matrix-multiplication} and Theorem \ref{thm:linear-maps-and-matrices-same-data}(i).
\end{enumerate}
\hlt{Structural Commentary}: Spectral reduction turns dimension induction into independent restricted families.\\
\hlt{Key Insight}: Commutants preserve all spectral subspaces.
\end{proof}
\end{theorem}

Figure \ref{fig:simultaneous-normal-diagonalization} illustrates the spectral-subspace reduction in Theorem \ref{thm:simultaneous-unitary-diagonalization-normal-family}.

\begin{illustrationfigure}[H]
  \centering
  \begin{tikzpicture}[
  >=Stealth,
  every node/.style={font=\small},
  block/.style={draw=gray!55, rounded corners=3pt, inner xsep=8pt, inner ysep=4.5pt, align=center},
  arrow/.style={->, line width=0.8pt, blue!70!black, shorten >=2pt, shorten <=2pt}
]
  \node[block] (family) at (0,0) {$\mathcal{F}$};
  \node[block] (split) at (2.46,0) {$A_0=\sum_j\mu_jP_j$};
  \node[block] (pieces) at (5.32,0) {$\bigoplus_j\operatorname{range}(P_j)$};
  \node[block] (basis) at (8.49,0) {common eigenbasis};

  \draw[arrow] (family.east) -- (split.west);
  \draw[arrow] (split.east) -- (pieces.west);
  \draw[arrow] (pieces.east) -- (basis.west);
\end{tikzpicture}

  \caption{The spectral ranges of $A_0=\sum\limits_j\mu_jP_j$ reduce the commuting family for dimension induction.}
  \label{fig:simultaneous-normal-diagonalization}
\end{illustrationfigure}

\begin{corollary} \label{cor:commuting-normal-sums-products-normal} \hlt{Sums and Products of Commuting Normal Operators}\\
Let $\mathcal{V}$ be a complex inner-product space, and let $S,T\in\LL(\mathcal{V})$ be commuting normal operators. Then $S+T$ and $ST$ are normal.
\begin{proof}
Fix a common orthonormal eigenbasis for $S$ and $T$ by Theorem \ref{thm:simultaneous-unitary-diagonalization-normal-family}.\\
The operators $S+T$ and $ST$ are diagonal in that basis and therefore normal by Theorem \ref{thm:complex-spectral-theorem}(iii),(i).\\
\hlt{Structural Commentary}: Simultaneous diagonalization reduces both constructions to scalar arithmetic.\\
\hlt{Key Insight}: Common diagonal coordinates are closed under sums and products.
\end{proof}
\end{corollary}

\begin{perspective} \label{per:real-normal-operators-planes} \hlt{Real Normal Operators Use Planes}\\
Theorem \ref{thm:real-operator-invariant-line-or-plane} replaces complex eigenlines by invariant real lines and planes.\\
For normal operators, orthogonal reduction turns a nonreal conjugate pair into a rotation-dilation block.\\
This viewpoint is the real substitute for diagonalization; it does not make every invariant plane irreducible.
\end{perspective}

\begin{proposition} \label{prop:two-dimensional-real-normal-block} \hlt{Two-Dimensional Real Normal Blocks}\\
Let $\mathcal{V}$ be a two-dimensional real inner-product space and let $T\in\LL(\mathcal{V})$.
\begin{enumerate}[label=(\roman*)]
\setlength{\itemsep}{0pt}
\item $T$ is normal and not self-adjoint if and only if every orthonormal-basis matrix has the form $\left(\begin{smallmatrix}a&-b\\b&a\end{smallmatrix}\right)$.\\
The corresponding parameters satisfy $a\in\mathbb{R}$ and $b\in\mathbb{R}\setminus\{0\}$.
\item If the equivalent conditions in (i) hold, then an orthonormal basis exists for which $b>0$.
\end{enumerate}
\begin{proof}
\begin{enumerate}[label=(\roman*)]
\setlength{\itemsep}{0pt}
\item[$(i)$] ($\Rightarrow$) Fix an arbitrary orthonormal basis and write $A=[T]=\left(\begin{smallmatrix}a&c\\b&d\end{smallmatrix}\right)$.\\
The identities $A^tA=AA^t$ and $A\neq A^t$ hold by Theorem \ref{thm:matrix-of-adjoint-is-conjugate-transpose} and Proposition \ref{prop:composition-and-matrix-multiplication}.\\
Equality of the $(1,1)$ entries is $b^2=c^2$. The relation $c\neq b$ holds by $A\neq A^t$, so $c=-b$ and $b\neq0$.\\
Equality of the off-diagonal entries is $-ab+bd=ab-bd$, so $d=a$.\\
The stated form holds in every orthonormal basis by arbitrariness of the basis.
\item[$(i)$] ($\Leftarrow$) Fix an orthonormal basis with $A=[T]=\left(\begin{smallmatrix}a&-b\\b&a\end{smallmatrix}\right)$ and $b\neq0$.\\
The identities $A^tA=AA^t$ and $A\neq A^t$ hold by direct multiplication and $b\neq0$.\\
Thus $T$ is normal and not self-adjoint by Theorem \ref{thm:matrix-of-adjoint-is-conjugate-transpose} and Proposition \ref{prop:composition-and-matrix-multiplication}.
\item[$(ii)$] Fix an orthonormal basis and its nonzero parameter $b$ by part (i).\\
If $b<0$, negate the second basis vector; the new parameter is $-b>0$. If $b>0$, retain the basis.
\end{enumerate}
\hlt{Structural Commentary}: Entrywise normality determines the block, followed by an orientation choice.\\
\hlt{Key Insight}: Nonsymmetry selects $c=-b$ from $c^2=b^2$.
\end{proof}
\end{proposition}

\begin{theorem} \label{thm:real-normal-operator-block-diagonal-form} \hlt{Real Normal Operator Block Form}\\
Let $\mathcal{V}$ be a real inner-product space and let $T\in\LL(\mathcal{V})$. For $a\in\mathbb{R}$ and $b>0$, set $N_{a,b}\defeq\left(\begin{smallmatrix}a&-b\\b&a\end{smallmatrix}\right)$.\\
Then $T$ is normal if and only if some orthonormal basis makes $[T]$ a direct sum of scalar blocks and blocks $N_{a,b}$.
\begin{proof}
\begin{enumerate}[label=\roman*.]
\setlength{\itemsep}{0pt}
\item[$(\Rightarrow)$] Induct on $n\defeq\dim(\mathcal{V})$. For $n=0$, the empty orthonormal basis has the required form.\\
For $n=1$, any unit vector produces a scalar block. Suppose $n\geq2$. Fix a $T$-invariant $\mathcal{U}$ with $\dim(\mathcal{U})\in\{1,2\}$ by Theorem \ref{thm:real-operator-invariant-line-or-plane}.\\
The subspace $\mathcal{U}^{\perp}$ is $T$-invariant, and both restrictions are normal, by Corollary \ref{cor:normal-operator-invariant-subspaces-reduce}(i),(iv).\\
If $\dim(\mathcal{U})=1$, a unit vector in $\mathcal{U}$ determines a scalar block.\\
If $\dim(\mathcal{U})=2$ and $T|_{\mathcal{U}}$ is self-adjoint, fix an orthonormal eigenbasis by Theorem \ref{thm:real-spectral-theorem}(i),(ii). Its two blocks are scalar.\\
If $T|_{\mathcal{U}}$ is not self-adjoint, fix an orthonormal basis with matrix $N_{a,b}$ by Proposition \ref{prop:two-dimensional-real-normal-block}(ii).\\
Fix the required basis on $\mathcal{U}^{\perp}$ by the induction hypothesis.\\
Their concatenation is an orthonormal basis of $\mathcal{V}=\mathcal{U}\oplus\mathcal{U}^{\perp}$ by Theorem \ref{thm:orthogonal-direct-sum-decomposition}.
\item[$(\Leftarrow)$] Fix an orthonormal basis with the stated block form and set $A\defeq[T]$.\\
Every allowed block $D$ satisfies $D^*D=DD^*$ by direct multiplication, so $A^*A=AA^*$ blockwise.\\
The identity $T^*T=TT^*$ holds by Theorem \ref{thm:matrix-of-adjoint-is-conjugate-transpose}, Proposition \ref{prop:composition-and-matrix-multiplication}, and Theorem \ref{thm:linear-maps-and-matrices-same-data}(i).\\
Thus $T$ is normal by its defining identity.
\end{enumerate}
\hlt{Structural Commentary}: Orthogonal reduction supports induction over invariant lines and planes.\\
\hlt{Key Insight}: Two-dimensional nonsymmetric blocks have rotation-dilation form.
\end{proof}
\end{theorem}

Figure \ref{fig:real-normal-block-structure} summarizes Theorem \ref{thm:real-normal-operator-block-diagonal-form}.

\begin{illustrationfigure}[H]
  \centering
  \begin{tikzpicture}[
  >=Stealth,
  every node/.style={font=\small},
  box/.style={draw=gray!55, rounded corners=3pt, inner xsep=8pt, inner ysep=6pt, align=center},
  arrow/.style={->, line width=0.85pt, blue!70!black}
]
  \node[box] (normal) at (0,0) {normal $T$};
  \node[box] (sum) at (3.6,0) {$\mathcal{V}=\bigoplus L_j\oplus\bigoplus P_k$};
  \node[box] (blocks) at (7.8,0) {$[T]=\operatorname{diag}(\lambda_j,R_k)$\\[-1pt]
  $R_k=\begin{bmatrix}a_k&-b_k\\ b_k&a_k\end{bmatrix}$};

  \draw[arrow] (normal.east) -- (sum.west);
  \draw[arrow] (sum.east) -- (blocks.west);
\end{tikzpicture}

  \caption{Orthogonal line/plane split into scalar and rotation blocks.}
  \label{fig:real-normal-block-structure}
\end{illustrationfigure}

\begin{theorem} \label{thm:complex-isometry-description} \hlt{Description of Complex Isometries}\\
Let $\mathcal{V}$ be a complex inner-product space, $n\defeq\dim(\mathcal{V})$, and $S\in\LL(\mathcal{V})$.\\
Then $S$ is an isometry if and only if $\mathcal{V}$ has an orthonormal eigenbasis $(\mathbf{e}_j)_{j=1}^{n}$ with eigenvalues $(\lambda_j)_{j=1}^{n}$ satisfying $(\forall j\in J_n)\,\abs{\lambda_j}=1$.
\begin{proof}
If $\mathcal{V}=\{\mathbf{0}\}$, both conditions hold for the unique operator and empty orthonormal basis. Assume $\mathcal{V}\neq\{\mathbf{0}\}$.
\begin{enumerate}[label=\roman*.]
\setlength{\itemsep}{0pt}
\item[$(\Rightarrow)$] Suppose $S$ is an isometry. Then $S^*S=I_{\mathcal{V}}$ by Theorem \ref{thm:finite-dimensional-isometry-iff-adjoint-product-identity}.\\
The map $S$ is injective by norm preservation and invertible by Theorem \ref{thm:finite-dimensional-operator-invertible-equivalences}(ii),(i). Thus $S^*=S^{-1}$ and $SS^*=I_{\mathcal{V}}$, so $S$ is normal.\\
Fix an orthonormal eigenbasis $S\mathbf{e}_j=\lambda_j\mathbf{e}_j$ by Theorem \ref{thm:complex-spectral-theorem}(i),(ii).\\
Then $(\forall j\in J_n)\,\abs{\lambda_j}=\norm{S\mathbf{e}_j}=1$ by norm preservation.
\item[$(\Leftarrow)$] Suppose $(\mathbf{e}_j)_{j=1}^{n}$ is an orthonormal eigenbasis with eigenvalues $(\lambda_j)_{j=1}^{n}$ satisfying $(\forall j\in J_n)\,\abs{\lambda_j}=1$.\\
For $\mathbf{v}=\sum\limits_{j=1}^{n}c_j\mathbf{e}_j$, $\norm{S\mathbf{v}}^2=\sum\limits_{j=1}^{n}\abs{\lambda_jc_j}^2=\sum\limits_{j=1}^{n}\abs{c_j}^2=\norm{\mathbf{v}}^2$ by Theorem \ref{thm:parseval-identity-finite-dimensional}(i). Hence $S$ is an isometry.
\end{enumerate}
\hlt{Structural Commentary}: A forward normality reduction pairs with a reverse Parseval calculation.\\
\hlt{Key Insight}: Isometry becomes coefficientwise unit modulus.
\end{proof}
\end{theorem}

\begin{example} \label{ex:plane-rotation-isometry} \hlt{Plane Rotations}\\
For $\theta\in\mathbb{R}$, set $R_{\theta}\defeq\left(\begin{smallmatrix}\cos\theta&-\sin\theta\\\sin\theta&\cos\theta\end{smallmatrix}\right)$. The identity $R_{\theta}^tR_{\theta}=I_2$ holds by direct multiplication.\\
Thus $R_{\theta}$ is an isometry of $\mathbb{R}^2$ by Theorem \ref{thm:matrix-of-adjoint-is-conjugate-transpose}, Proposition \ref{prop:composition-and-matrix-multiplication}, and Theorem \ref{thm:finite-dimensional-isometry-iff-adjoint-product-identity}.\\
If $\theta\notin\pi\mathbb{Z}$, its characteristic polynomial $x^2-2(\cos\theta)x+1$ has discriminant $-4\sin^2\theta<0$.\\
Hence $R_{\theta}$ has no real eigenvalue, which motivates two-dimensional blocks in the real isometry form.
\end{example}

\begin{theorem} \label{thm:real-isometry-block-diagonal-form} \hlt{Real Isometry Block Form}\\
Let $\mathcal{V}$ be a real inner-product space and let $S\in\LL(\mathcal{V})$. For $\theta\in(0,\pi)$, set $R_{\theta}\defeq\left(\begin{smallmatrix}\cos\theta&-\sin\theta\\\sin\theta&\cos\theta\end{smallmatrix}\right)$.\\
Then $S$ is an isometry if and only if some orthonormal basis makes $[S]$ a direct sum of $[1]$, $[-1]$, and blocks $R_{\theta}$.
\begin{proof}
If $\mathcal{V}=\{\mathbf{0}\}$, both conditions hold for the unique operator and empty orthonormal basis. Assume $\dim(\mathcal{V})\geq1$.
\begin{enumerate}[label=\roman*.]
\setlength{\itemsep}{0pt}
\item[$(\Rightarrow)$] Suppose $S$ is an isometry. Then $S^*S=I_{\mathcal{V}}$ by Theorem \ref{thm:finite-dimensional-isometry-iff-adjoint-product-identity}.\\
The map $S$ is injective by norm preservation and invertible by Theorem \ref{thm:finite-dimensional-operator-invertible-equivalences}(ii),(i).\\
Hence $S^*=S^{-1}$, so $SS^*=I_{\mathcal{V}}$ and $S$ is normal. Fix a block basis by Theorem \ref{thm:real-normal-operator-block-diagonal-form}.\\
In this basis, $[S]^t[S]=I$ by Theorem \ref{thm:matrix-of-adjoint-is-conjugate-transpose} and Proposition \ref{prop:composition-and-matrix-multiplication}. Every scalar block $[\lambda]$ satisfies $\lambda^2=1$, so $\lambda=\pm1$.\\
Every two-dimensional block $\left(\begin{smallmatrix}a&-b\\b&a\end{smallmatrix}\right)$ satisfies $a^2+b^2=1$ with $b>0$.\\
Hence a unique $\theta\in(0,\pi)$ satisfies $a=\cos\theta$ and $b=\sin\theta$, so the block equals $R_{\theta}$.
\item[$(\Leftarrow)$] Fix a basis with the stated block form. Every block $D$ satisfies $D^tD=I$.\\
Thus $[S]^t[S]=I$, and $S^*S=I_{\mathcal{V}}$ by Theorem \ref{thm:matrix-of-adjoint-is-conjugate-transpose}, Proposition \ref{prop:composition-and-matrix-multiplication}, and Theorem \ref{thm:linear-maps-and-matrices-same-data}(i).\\
Hence $S$ is an isometry by Theorem \ref{thm:finite-dimensional-isometry-iff-adjoint-product-identity}.
\end{enumerate}
\hlt{Structural Commentary}: The real normal form is filtered blockwise by the isometry identity.\\
\hlt{Key Insight}: Unit-length normal blocks are signs or rotations.
\end{proof}
\end{theorem}

Figure \ref{fig:real-isometry-blocks} illustrates the three block types in Theorem \ref{thm:real-isometry-block-diagonal-form}.

\begin{illustrationfigure}[H]
  \centering
  \begin{tikzpicture}[
  >=Stealth,
  every node/.style={font=\small},
  box/.style={draw=gray!55, rounded corners=3pt, inner xsep=8pt, inner ysep=6pt, align=center},
  arrow/.style={->, line width=0.85pt, blue!70!black}
]
  \node[box] (iso) at (0,0) {isometry $S$};
  \node[box] (blocks) at (3.3,0) {orthogonal sum\\$[1]$, $[-1]$, $R_{\theta}$};
  \node[box] (rot) at (7.2,0) {$R_{\theta}=\begin{bmatrix}\cos\theta&-\sin\theta\\ \sin\theta&\cos\theta\end{bmatrix}$};

  \draw[arrow] (iso) -- (blocks);
  \draw[arrow] (blocks) -- (rot);
\end{tikzpicture}

  \caption{Normal form: fixed/flipped lines and rotation planes.}
  \label{fig:real-isometry-blocks}
\end{illustrationfigure}
